% manuscript/drafts/related_work.tex
%
% Session L4, Step 13 draft. Organised by the theory areas of the Step 2 query
% matrix, so coverage is auditable against docs/literature/QUERY_MATRIX.md.
%
% Precise detail is drawn from docs/literature/ANCHOR_NOTES.md, which records page
% and equation pointers for every paper actually read, and marks which papers were
% verified by identifier only. Nothing page-precise is claimed from an unread paper.

\section{Related work}
\label{sec:related}

\subsection{Rossby-wave dispersion on a sphere}

The westward propagation of large-scale atmospheric disturbances as a consequence
of the latitudinal variation of the Coriolis parameter is due to
\citet{rossby1939}, who obtained it on a tangent plane, and to Haurwitz, who
extended it first to disturbances of finite lateral extent and then to the sphere,
where the closed form $c_{\mathrm{ang}} = -2\Omega/[n(n+1)]$ for a spherical
harmonic of degree $n$ replaces the plane-wave result $c = -\beta/k^2$. The
spherical solutions are exact steadily propagating solutions of the nonlinear
nondivergent barotropic vorticity equation, which is what made them attractive as
model test cases.

That attractiveness has a limit which this paper takes seriously.
\citet{thuburn2000} integrated Rossby--Haurwitz waves in four independent
high-resolution shallow-water models and established two things that bear directly
on any campaign measuring their phase speed. First, the shallow-water
Rossby--Haurwitz wave --- unlike its nondivergent counterpart --- generates
small-scale structure and therefore sustains a potential-enstrophy cascade, so its
numerical solution depends on the dissipation coefficient, and the appropriate
coefficient depends on the flow. Second, and contrary to the belief on which the
standard test case was built, the zonal-wavenumber-4 wave is dynamically unstable
and breaks down when perturbed, with exponential growth on an e-folding time of
order a day. \citet{lynch2009} approaches the same question from the resonant-triad
direction.

The stability of Rossby--Haurwitz waves has since been treated with increasing
rigour: \citet{skiba2004} computed the structure and growth rates of unstable modes
to the Rossby--Haurwitz wave, \citet{skiba2008,skiba2024} developed the linear and
nonlinear instability theory on the sphere, and \citet{constantin2022} gave a
rigorous treatment of stratospheric planetary flows from the Euler-equation
perspective.

\subsection{The divergent correction: Laplace's tidal equations and Hough modes}

Retaining the free surface converts the nondivergent vorticity problem into
Laplace's tidal equations, whose eigenfunctions on a sphere were computed by
\citet{longuethiggins1968} as a function of Lamb's parameter, and recast in vector
harmonic form by \citet{swarztrauber1985}. \citet{kasahara1976} reviews the
construction and states the branch structure explicitly: the horizontal parts of
the normal modes are Hough harmonics, carrying zonal velocity, meridional velocity
and geopotential height together, and they fall into three families of distinct
frequency --- eastward- and westward-propagating gravity waves, and
westward-propagating rotational waves of Rossby--Haurwitz type.

Two remarks on how this paper uses that body of work. The dependence of the
rotational-branch frequencies on Lamb's parameter --- equivalently, the departure
from the nondivergent $-2\Omega/[n(n+1)]$ as the deformation radius shrinks
relative to the planetary radius --- is classical, and is not claimed here as a new
result. What this paper does with it is narrower: it uses the closed-form
$\varepsilon \to 0$ limit as the validation target for its own eigenvalue solver,
so that the analytic reference is one the study derives and checks internally
rather than a table transcribed from elsewhere. Of the three classical sources,
only \citet{kasahara1976} was available to us in full; \citet{longuethiggins1968}
and \citet{swarztrauber1985} are cited on the strength of their bibliographic
records, and no page-precise claim is made from either.

\subsection{Barotropic instability of a zonal jet}

\citet{kuo1949} established the necessary condition for instability of a
two-dimensional nondivergent barotropic flow: a growing normal mode requires a
critical point at which the gradient of absolute vorticity $\beta - \bar u_{yy}$
vanishes, so that the background gradient changes sign somewhere in the domain. He
is explicit about the limits of that result, noting that dynamic stability analysis
indicates only the possibility of development and that the mechanism of the
disturbance must be studied separately.

\citet{bretherton1966} generalised the constraint and supplied its physical
reading. His integrated eddy potential-vorticity flux relation shows that a growing
normal mode requires the background potential-vorticity gradient to take both
signs, recovering for a wider class of flows the result Charney and Stern had
proved for a special case; and his observation that a discontinuity in the
background field is mathematically equivalent to a delta function in potential
vorticity is what allows the interface idealisation used later in this paper to be
exact for piecewise-constant $Q$ rather than merely illustrative. His main result
concerns the critical layer: a growing mode of a general zonal flow must have one,
and it carries a down-gradient potential-vorticity flux that does not vanish as the
growth rate does.

For the divergent shallow-water system the stability theory is different, and
\citet{ripa1983} is the reference: general stability conditions for zonal flows in
a one-layer model on the beta-plane or the sphere. \citet{hayashi1987} treat
stable and unstable shear modes of rotating parallel flows in shallow water, and
\citet{knessl1992} the stability of rotating shear flows in the same system.
This paper's own stability analysis is nondivergent, and
\S\ref{sec:discussion} states the consequence rather than leaving it implicit.

\subsection{Counter-propagating Rossby waves}

\citet{bretherton1966} showed that development in a flow with two distinct
potential-vorticity gradients is the interaction of two Rossby waves propagating on
them; \citet{hoskins1985} developed the potential-vorticity-inversion framework
that makes the ``wind field attributable to'' each wave a well-defined object, and
named the counter-propagating picture. \citet{heifetz1999} gave the quantitative
form for Rayleigh's constant-shear strip, tracking the two edge waves' amplitudes
and phases explicitly and showing that the growing normal mode is the phase-locked
configuration: growth occurs when the northern wave lies less than half a
wavelength west of the southern one. Their dispersion relation and the constants
that follow from it --- a short-wave cutoff at nondimensional wavenumber
$K_{\mathrm c} \approx 1.28$, a most-unstable $K_{\max} \approx 0.8$ corresponding
to a wavelength about eight times the vorticity-strip width, and a peak growth rate
of about twenty per cent of the shear --- are the external check on the
two-interface reduction used in this paper.

Their framing should be adopted along with their algebra. \citet{heifetz1999} offer
the counter-propagating description because, in their words, it provides a useful
pedagogical framework, having found the normal-mode explanation --- that the waves
grow because a complex phase speed is required for the boundary conditions to be
satisfiable --- unsatisfying. The one-interface/two-interface identification in this
paper is exposition in the same spirit, not a new mechanism.

\subsection{The Rhines scale and jet formation}

\citet{rhines1975} showed that on a beta-plane the upscale cascade of
two-dimensional turbulence is arrested near a wavenumber
$k_\beta = (\beta/2U)^{1/2}$, set only by the root-mean-square particle speed and
$\beta$, and that the cascade generates zonal flow by itself. Two features of that
result constrain how it may be used. The arrest is anisotropic --- the restoring
term acts only on meridional motion, so energy accumulates in zonally elongated
structures --- and the alternating-jet end state is reached in the homogeneous case,
Rhines noting that when energy is spatially intermittent the cascade may be halted
instead by the spreading of energy about space, with the zonal end state probably
not achieved. \citet{vallis1993} reach the same $O(\sqrt{U/\beta})$ transition
scale from the weakly nonlinear interaction of Rossby waves, and emphasise that
differential rotation produces anisotropy directly rather than merely inhibiting
the cascade; they also note that the particular form of the scale is not crucial to
their argument, which is the basis for treating its numerical factor as
convention-dependent.

This paper uses the Rhines scale for one purpose only: because
$\beta \propto \Omega$, sweeping the rotation rate sweeps $L_{\mathrm R}$, which
reframes the rotation sweep as a question about planetary regime. It does not treat
$L_{\mathrm R}$ as a jet-spacing law, a stronger claim than either
\citet{rhines1975} or \citet{vallis1993} supports.

\subsection{Benchmarks, reference solutions and the numerical framework}

\citet{williamson1992} proposed the standard test set for shallow-water
approximations in spherical geometry that the field still uses, with pseudospectral
reference solutions supplied by \citet{jakobchien1995}. Its limitations are
documented by its users: \citet{galewsky2004} record that widespread use revealed
problems severely limiting the suite's practical utility --- the early cases being
idealised to the point of being unrepresentative, the later ones carrying
unexpected subtleties --- and propose in its place a barotropically unstable
mid-latitude jet, analytically specified, whose evolution exercises both the fast
gravity-wave and slow balanced parts of the shallow-water system. That test case
has since become a de facto acceptance test: the great majority of the work citing
it is discretisation development --- finite-volume, discontinuous-Galerkin,
icosahedral and cubed-sphere, radial-basis-function and semi-Lagrangian schemes ---
rather than analysis of the jet's own dynamics. \citet{lauter2005} supply an
unsteady analytic solution of the spherical shallow-water equations, which this
paper adopts as an additional verification case complementing the steady one.

The solver is Dedalus \citep{burns2020}, a spectral framework for partial
differential equations, using the spherical basis and tensor calculus of
\citet{vasil2019}. Both are used as documented; no numerical development is claimed
here.

\subsection{Observations}

The comparison against observation uses two independent reanalyses, ERA5
\citep{hersbach2020} and the NCEP/NCAR reanalysis \citep{kalnay1996}, so that the
spread between them can be reported as observational uncertainty rather than
concealed by choosing one. The diagnostic is a space--time spectral decomposition
of geopotential height, which separates eastward- and westward-propagating branches
explicitly --- necessary here because the observed extratropical troposphere is
dominated by eastward-moving baroclinic disturbances that a single-layer barotropic
model does not represent, and a simple temporal high-pass filter would retain them
mixed in with any westward barotropic signal.

Two points of honesty about this comparison. Correcting an observed
ground-relative phase speed for advection by the mean flow before comparing it with
an intrinsic theoretical prediction is standard practice, not a contribution;
omitting it would be an error, and performing it earns no credit. And the
comparison is bounded to scales rather than trajectories: agreement is sought in
dominant zonal wavenumber and in the order of magnitude of drift speed, with
disagreement in detail expected from the baroclinicity, vertical structure and
forcing that the model omits.
